\documentclass[10pt, fleqn]{article}
\usepackage[a4paper,margin=0.5in]{geometry}
\usepackage{multicol}
\usepackage{amsmath, amssymb}
\usepackage{titlesec}
\usepackage{float}
\usepackage{amsfonts, bm, graphicx}
\usepackage{enumitem}
\usepackage[font=small, skip=4pt]{caption}

\setlength{\abovedisplayskip}{6pt}
\setlength{\belowdisplayskip}{6pt}
\setlength{\abovedisplayshortskip}{4pt}
\setlength{\belowdisplayshortskip}{4pt}

\setlength{\textfloatsep}{8pt plus 2pt minus 2pt}
\setlength{\intextsep}{6pt plus 2pt minus 2pt}
\setlength{\floatsep}{8pt plus 2pt minus 2pt}

\setlength{\mathindent}{0pt}

\titleformat{\section}{\large\bfseries}{}{0em}{}
\titleformat{\subsection}{\normalsize\bfseries}{}{0em}{}
% Uppercase Volume with thicker, higher line
\newcommand{\vol}{\mathord{\text{\ooalign{\hidewidth $V$\hidewidth\cr\kern-0.5pt\raisebox{0.8ex}{\rule[0pt]{1em}{0.5pt}}}}}}

% Lowercase Volume with manual horizontal shift
\newcommand{\svol}[1][-0.15em]{% optional argument is horizontal shift
  \mathord{\text{\ooalign{
    $v$\cr
    \hspace{#1}\rule[0.6ex]{0.8em}{0.5pt} % 0.6ex vertical raise, 0.8em length, 0.5pt thickness
  }}}%
}
\begin{document}

\begin{center}
    \Large \textbf{Thermodynamics Formula Sheet} \\
    \normalsize Chapter 2
\end{center}

\begin{multicols}{2}

\section*{Pressure}
\textbf{Constants}
\[\begin{aligned}
1 \text{ atm} &= 101.325 \text{ kPa} \\
&= 1.01325 \text{ bar} \\
&= 760 \text{ Torr} \\
&= 760 \text{ mmHg} \\
&= 29.92 \text{ inHg} \\
&= 14.696 \text{ psi}\\
&= 10.33 \text{ m H}_2\text{O (at 4$^\circ$C)} \\
\end{aligned}\]
\textbf{Pressure Conversions}
\[1\,\text{Pa} = 1\,\text{N/m}^2\]
\[1\, \text{ bar} = 10^{5}\, \text{ Pa}\]
\[1\, \text{ psi} = 6.89476 \times 10^{3}\, \text{ Pa}\]
\[1\, \text{ Torr} = 133.322\, \text{ Pa}\]
\[1\, \text{ mmHg} = 133.322\, \text{ Pa}\]
\[1\, \text{ inHg} = 3.38639 \times 10^{3}\, \text{ Pa}\]
\textbf{Gauge vs Absolute Pressure}
\[P_{\text{gauge}} = P_{\text{absolute}} - P_{\text{atmospheric}}\]
\[P_{\text{vacuum}} = P_{\text{atmospheric}} - P_{\text{absolute}}\]
\underline{Note:} Generally, gauge pressures already account for atmospheric pressure, and thus reads zero when open to the atmosphere.\\
\textbf{Pressure Formulas}
\[P=F/A\]
\[P = \rho g h\quad \text{(hydrostatic pressure)}\]
\section*{Temperature}
\textbf{Temperature Conversions}
\begin{itemize}[label=\textbullet]
    \item Absolute temperature conversions
    \begin{itemize}[label=\textopenbullet]
        \item $T_{K} = T_{^\circ C} + 273.15$
        \item $T_{^\circ R} = T_{^\circ F} + 459.67$
        \item $T_{^\circ F} = \frac{9}{5}T_{^\circ C} + 32$
        \item $T_{^\circ R} = \frac{9}{5}T_{K}$
    \end{itemize}
    \item Temperature difference conversions
    \begin{itemize}[label=\textopenbullet]
        \item $\Delta K = \Delta ^\circ C$
        \item $\Delta ^\circ R = \Delta ^\circ F$
        \item $\Delta ^\circ F = \frac{9}{5}\Delta ^\circ C$
        \item $\Delta ^\circ R = \frac{9}{5}\Delta K$
    \end{itemize}
\end{itemize}
\section*{SI Prefixes}
\makeatletter
\@fleqnfalse
\makeatother
\[\begin{aligned}
\text{femto (f)}  &= 10^{-15} \\
\text{pico (p)}   &= 10^{-12} \\
\text{nano (n)}   &= 10^{-9} \\
\text{micro ($\mu$)}  &= 10^{-6} \\
\text{milli (m)}  &= 10^{-3} \\
\text{centi (c)}  &= 10^{-2} \\
\text{deci (d)}   &= 10^{-1} \\
\text{deca (da)}  &= 10^{1} \\
\text{hecto (h)}  &= 10^{2} \\
\text{kilo (k)}   &= 10^{3} \\
\text{mega (M)}   &= 10^{6} \\
\text{giga (G)}   &= 10^{9} \\
\text{tera (T)}   &= 10^{12} \\
\text{peta (P)}   &= 10^{15} \\
\end{aligned}\]
\makeatletter
\@fleqntrue
\makeatother
\underline{Note:} $1\text{Angstrom}\,(\text{\AA}) = 10^{-10}$ m
\section*{Density and Specific Gravity}
\[\rho = \frac{m}{\vol}\quad \text{(density)}\]
\[\svol = \frac{\vol}{m} = \frac{1}{\rho} \quad \text{(specific volume)}\]
\textbf{Specific gravity} is defined as a relative density compared to water (at 4$^\circ$C) where $\rho_{water} = 1000\,\text{kg/m}^3$:
\[\boxed{SG = \frac{\rho_{fluid}}{\rho_{water}} \implies \rho_{fluid} = SG \cdot \rho_{water}}\]
\textbf{Specific Weight}
\[\gamma_s = \rho g \quad \text{(specific weight)}\]
Where hydrostatic pressure can be written as $P = \gamma_s h$\\\\
\textbf{Slugs and lbm}
\[\left[\frac{1\,\text{lbf}}{32.174\,\frac{\text{lbm}\cdot\text{ft}}{\text{s}^2}}\right],\quad
\left[\frac{1\,\text{slug}}{32.174\,\text{lb,}}\right], \quad
\left[\frac{1\,\text{lbf}}{1\,\frac{\text{slug}\cdot\text{ft}}{\text{s}^2}}\right]\]
\textbf{Energy Units}\\
1 Btu = 778.17 ft$\cdot$lbf = 1055.06 J\\\\
\textbf{Work Units}\\
1 hp = 550 ft$\cdot$lbf/s = 745.7 W
\end{multicols}\newpage
\makeatletter
\@fleqnfalse
\makeatother
\section*{Power}
Power ($\dot{W}$) is the rate of work done, and can be calculated for a control volume (open system with mass flow) using the first law of thermodynamics:
\[\boxed{\frac{dE}{dt} = \dot{Q}-\dot{W}+\sum \dot{m_i}\left[h_i+\frac{1}{2}v_i^2+g z_i\right]-\sum \dot{m_e}\left[h_e+\frac{1}{2}v_e^2+g z_e\right]}\]
Where:
\begin{itemize}
\item $dE/dt$ is the rate of change of the total energy of the system over time
\item $\dot{Q}$ is the rate of heat transfer into the system from the surroundings
\item $\dot{W}$ is the rate of work done by the system on its surroundings (power output)
\item The summations represent the total energy entering and leaving the system due to mass flow, where $h$ is the specific enthalpy, $v/2$ is the specific kinetic energy, and $gz$ is the specific potential energy at the inlet ($i$) and outlet ($e$) of the control volume.
\end{itemize}
For a steady state process ($dE/dt = 0$) with no heat transfer ($\dot{Q}=0$), the equation becomes
\[\dot{W} = \dot{m} e_{\text{mech}}\]
\[e_{\text{mech}} = \frac{P}{\rho} +\frac{v^2}{2}+gz \quad \text{(P is pressure)}\]
Where $e_{\text{mech}}$ is the mechanical energy per unit mass, and $\dot{m}$ is the mass flow rate. 
\[\boxed{\dot{m}= \rho v A}\]
\[\boxed{\dot{m}= \rho\, \dot{\vol}}\]
\underline{Note:} $\dot{\vol}$ is the volumetric flow rate, and $A$ is the cross-sectional area of the flow.\\
From physics, power is also defined as 
\[\dot{W} = F \cdot v\]
Where $F$ is the force applied to an object and $v$ is the velocity of the object.\\\\
\textbf{Fluid Flow}
\[\boxed{\dot{W} = \dot{\vol}(P_e - P_i)}\]
\textbf{Spring}
\[\boxed{\dot{W}_{\text{spring}} = \frac{1}{2} k (x_f^2 - x_i^2)}\]
\section*{Heat Transfer by Conduction}
\[\boxed{\dot{Q}_{\text{cond}} = kA\frac{\Delta T}{\Delta x}}\]
Where: $\Delta T = T_i - T_e$, $A$ is the area of the window, $k$ is the thermal conductivity of the material, and $\Delta x$ is the thickness of the material. The amount of heat transferred is then given by
\[\boxed{Q_{\text{cond}} = \dot{Q}_{\text{cond}} \cdot t}\]
Where $t$ is the time duration of heat transfer.


\end{document}
